\documentclass[preprint,12pt]{elsarticle}

% ---------- Packages ----------
\usepackage{amsmath,amssymb,amsthm,mathtools}
\usepackage{enumitem}

% ---------- Theorem environments ----------
\newtheorem{proposition}{Proposition}
\newtheorem{corollary}[proposition]{Corollary}

\theoremstyle{definition}
\newtheorem{definition}[proposition]{Definition}
\newtheorem{remark}[proposition]{Remark}

% ---------- Commands ----------
\newcommand{\Lmu}{\mathcal{L}_{\mathrm{BA},\mu}}
\newcommand{\Tfa}{T_{\mathrm{fa}}}
\newcommand{\Q}{\mathbb{Q}}
\newcommand{\N}{\mathbb{N}}

\journal{Annals of Pure and Applied Logic}

\begin{document}

\begin{frontmatter}

\title{Countable Additivity is Not First-Order Axiomatizable}

\author{Patrick S.\ Mahon}
\ead{pmahon@sfu.ca}
\address{Research Computing Group, Simon Fraser University, Burnaby, BC, Canada}

\begin{abstract}
In the first-order language of Boolean algebras equipped with a normalized finitely
additive charge, no first-order theory characterizes those models whose charge extends
to a $\sigma$-additive measure on the generated $\sigma$-algebra.  Equivalently,
countable additivity is not first-order axiomatizable among finitely additive probability
Boolean algebras.  The proof is a single ultraproduct argument: Dirac charges on the
finite-cofinite algebra are individually $\sigma$-additive, but their ultraproduct is
the finite-cofinite charge, which fails $\sigma$-additivity --- contradicting
\L{}o\'{s}'s theorem if any first-order axiomatization existed.
\end{abstract}

\begin{keyword}
countable additivity \sep first-order axiomatizability \sep ultraproduct \sep
Boolean algebra \sep finitely additive measure \sep \L{}o\'s's theorem

\MSC[2020] 03C20 \sep 28A60 \sep 60A10
\end{keyword}

\end{frontmatter}

A property of structures is first-order axiomatizable in a language $L$ if there exists
a set of $L$-sentences whose models are exactly the structures satisfying it.
\L{}o\'{s}'s theorem gives the key diagnostic: if a property fails in an ultraproduct
of structures that individually satisfy it, no first-order axiomatization exists.

Countable additivity distinguishes $\sigma$-additive measures from finitely additive
charges.  That it cannot be expressed by a single first-order sentence is clear: the
condition quantifies over countable sequences, unavailable in first-order logic.  Its
non-first-order character is treated as background knowledge in the probability logic
literature~\cite{Hoover1978,Keisler1985,FaginHalpernMegiddo1990}, but the sharper
question --- whether countable additivity admits a first-order \emph{axiomatization}
by any set of sentences, not merely a single formula --- does not appear to have been
settled by an explicit proof.

We settle it negatively.

\medskip

Let $\Lmu$ be the one-sorted first-order language with:
\begin{itemize}
  \item Boolean operations $\wedge, \vee, \neg, 0, 1$;
  \item for each rational $q \in [0,1] \cap \Q$, unary relation symbols
        $M_{\leq q}(x)$ and $M_{\geq q}(x)$, intended to mean
        $\mu(x) \leq q$ and $\mu(x) \geq q$ respectively.
\end{itemize}

A structure for $\Lmu$ is a Boolean algebra $B$ together with a $[0,1]$-valued finitely
additive probability $\mu$, encoded through the truth of the predicates $M_{\leq q}$
and $M_{\geq q}$.  Using rational comparison predicates rather than a function symbol
into $[0,1]$ avoids introducing a second numeric sort.

\begin{definition}[First-order theory of finitely additive probabilities]
Let $\Tfa$ be the $\Lmu$-theory axiomatizing:
\begin{enumerate}[label=(\roman*)]
  \item $B$ is a Boolean algebra;
  \item normalization: $\mu(0) = 0$ and $\mu(1) = 1$;
  \item monotonicity: $x \leq y \Rightarrow \mu(x) \leq \mu(y)$;
  \item finite additivity: for all rationals $r, s$ with $r + s \leq 1$,
        \[
          x \wedge y = 0 \wedge M_{\leq r}(x) \wedge M_{\leq s}(y)
          \;\Rightarrow\; M_{\leq r+s}(x \vee y),
        \]
        and similarly for lower bounds.
\end{enumerate}
\end{definition}

\noindent
The class of Boolean algebras with normalized finitely additive probability is
first-order axiomatizable in $\Lmu$.  Countable additivity, by contrast, requires
quantification over an arbitrary countable sequence $(a_n)_{n \in \N}$ --- not
available in first-order logic --- and so is not a sentence of $\Lmu$.

\medskip

\begin{proposition}[$\sigma$-additive extension is not first-order axiomatizable]
\label{prop:not-fo}
There is no first-order $\Lmu$-theory $T \supseteq \Tfa$ such that, for every
$(B, \mu) \models \Tfa$,
\[
  (B, \mu) \models T \quad\Longleftrightarrow\quad
  \mu \text{ extends to a $\sigma$-additive measure on } \sigma(B).
\]
Equivalently, countable additivity is not first-order axiomatizable among
finitely additive probability Boolean algebras.
\end{proposition}

\begin{proof}
Let $E$ be the Boolean algebra of finite-cofinite subsets of $\Q$, fix an
enumeration $q : \N \to \Q$, and let $\delta_n$ be the Dirac charge at
$q(n)$.  Each $(E, \delta_n)$ is a model of $\Tfa$ in which $\delta_n$ is
$\sigma$-additive.  Suppose for contradiction that a theory $T$ as described
exists; then each $(E, \delta_n) \models T$.

Let $\mathcal{U}$ be a nonprincipal ultrafilter on $\N$ extending the cofinite
filter.  By \L{}o\'{s}'s theorem~\cite{Los1955}, the ultraproduct
$\prod_\mathcal{U}(E,\delta_n)$ also models $T$.  The charge induced on the
diagonal copy of $E$ is
\[
  \ell(A) = \lim_\mathcal{U} \delta_n(A) = \lim_\mathcal{U} \mathbf{1}_{q(n) \in A}.
\]
If $A$ is finite, $q(n) \in A$ for only finitely many $n$, so $\ell(A) = 0$.
If $A$ is cofinite, $q(n) \in A$ for cofinitely many $n$, so $\ell(A) = 1$
since $\mathcal{U}$ extends the cofinite filter.  Thus $\ell$ is the
finite-cofinite charge.

The sequence $A_k = \Q \setminus \{q(0), \ldots, q(k)\}$ satisfies
$A_0 \supseteq A_1 \supseteq \cdots$ and $\bigcap_k A_k = \varnothing$, yet
$\ell(A_k) = 1$ for all $k$.  Hence the charge on the diagonal copy of $E$
fails $\sigma$-additivity and does not extend to a $\sigma$-additive measure
on $\sigma(E)$.  Consequently the ultraproduct cannot satisfy any theory $T$
characterizing $\sigma$-additive extension, contradicting \L{}o\'{s}'s theorem
applied to the assumption that each $(E, \delta_n) \models T$.
\end{proof}

\begin{remark}
The ultraproduct $\prod_\mathcal{U}(E,\delta_n)$ is an $\Lmu$-structure in its own
right; $\ell$ is the charge induced on the diagonal copy of $E$ inside it.
Identifying the ultraproduct with $(E,\ell)$ is harmless for the contradiction,
since the failure of $\sigma$-additivity is witnessed on the diagonal copy.
\end{remark}

\medskip

The proof produces one purely finitely additive charge as an ultraproduct of
$\sigma$-additive ones; that single witness suffices for the non-axiomatizability
conclusion.  A stronger question is:

\begin{quote}
\itshape
Which purely finitely additive charges on a Boolean algebra arise as pointwise
ultralimits of $\sigma$-additive probabilities on the same algebra?
\end{quote}

\noindent On a $\sigma$-algebra the answer is none: any pointwise ultralimit of
$\sigma$-additive probabilities is itself $\sigma$-additive, by Nikodym-type
convergence~\cite{DunfordSchwartz1958}, IV.9.8.  On a general Boolean algebra the
situation is more structured, and the problem is settled in several natural regimes.
The remaining open case reduces to whether there exists a non-$\sigma$-complete
non-atomic Boolean algebra admitting no $\sigma$-additive probability.

\bibliographystyle{elsarticle-num}
\bibliography{references}

\end{document}
